\documentclass[11pt]{article}

\usepackage[margin=1in]{geometry}
\usepackage{amsmath,amssymb}
\usepackage{mathtools}
\usepackage{booktabs}
\usepackage{hyperref}

\newcommand{\dg}{^{\dagger}}
\newcommand{\ket}[1]{|#1\rangle}
\newcommand{\braket}[1]{\langle #1\rangle}
\newcommand{\Tr}{\mathrm{Tr}}

\title{Why the direct-minimization embedding state is an excited eigenstate\\
\large (ghost Gutzwiller for the $t_c$-Cederbaum model, $B_c=3$)}
\author{}
\date{\today}

\begin{document}
\emergencystretch=4em
\maketitle

\noindent This note explains a numerical observation made when the solution of the direct minimization
(\texttt{not\_imp\_gga\_cederbaum.py}) was used as the starting point of the embedding fixed-point loop
(\texttt{gga\_cederbaum.py}) for the $t_c$-Cederbaum model of \texttt{cederbaum.tex}: the core embedding state $\Phi_c$ of the direct
minimization is an \emph{excited} eigenstate of the embedding Hamiltonian, which the loop (that takes the lowest state) cannot follow.
Notation is that of \texttt{cederbaum.tex}.

%======================================================================
\section{Statement}

``Excited'' does \emph{not} refer to the physical system: the direct-minimization point is the lowest-energy solution of the
constrained problem (its energy agrees with the exact ground-state energy to $\sim10^{-7}$). It refers to the embedding Hamiltonian
$H^c_{\rm imp}$ built from the multipliers recovered at that point: $\Phi_c$ is an exact eigenstate of $H^c_{\rm imp}$, but not the lowest one.

%======================================================================
\section{The two procedures and what each requires}

Both stationarize the same Lagrangian
\begin{equation}
\mathcal L=E+\sum_I2\,\Tr\big[\lambda_I\big({}^I\Delta(\Phi_I)-\rho^{II}\big)\big]-\sum_IE^c_I\big(\braket{\Phi_I|\Phi_I}-1\big),
\end{equation}
with $E$ of Eq.~(E) of \texttt{cederbaum.tex} and multipliers $\lambda_I$ (density constraint) and $E^c_I$ (normalization).

\paragraph{Direct minimization} minimizes $E$ over $\Phi_I$ and $\psi_{qp}$ subject to the constraints. At a constrained minimum, $\mathcal L$ is
stationary in $\Phi_I$, i.e.
\begin{equation}
H^I_{\rm imp}\,\Phi_I=E^c_I\,\Phi_I ,
\label{eq:eig}
\end{equation}
where, for the core,
\begin{equation}
H^c_{\rm imp}=H^{loc}_c+\sum_{a\sigma}h_{a}\big(d\dg_{a\sigma}c_\sigma+\mathrm{h.c.}\big)+\sum_{ab\sigma}\lambda^c_{ab}\,d\dg_{a\sigma}d_{b\sigma}
+n_c\Big[\alpha+2\Tr\Xi-\sum_{ab\sigma}\Xi_{ab}\,d\dg_{a\sigma}d_{b\sigma}\Big],\qquad\lambda^c=\Lambda-\lambda ,
\end{equation}
(with $h_a$, $\Lambda$, $\alpha$, $\Xi$ the derivatives of $E$ given in Sec.~Direct minimization of \texttt{cederbaum.tex}). Equation~\eqref{eq:eig} says that
$\Phi_I$ is \emph{an eigenstate} of $H^I_{\rm imp}$; it says nothing about which one.

\paragraph{The impurity loop} \emph{assumes} that $\Phi_I$ is the \emph{lowest} eigenstate of $H^I_{\rm imp}$ at the current multipliers (it diagonalizes
$H^I_{\rm imp}$ and takes the ground state). This is the saddle-point reading of $\mathcal L$: minimize over $\Phi_I$ at fixed multipliers,
then adjust the multipliers to satisfy the constraints. It is equivalent to the constrained minimum when the problem has no duality gap
(e.g.\ if it were convex).

%======================================================================
\section{Numerical check}

At $t_c=0.5$, $B_c=3$ ($\Delta=1$, $E_c=10$), for the two couplings $V/g=1/\sqrt2=0.707$ and $\sqrt2=1.414$ ($\rho/\Omega=V/(\sqrt2g)=0.5$ and $1.0$, the ratios of Kockl\"auner and Golze;
see \texttt{cederbaum.tex}), with the multipliers recovered from the direct-minimization point (shift $\varepsilon=10^{-9}$;
\texttt{excited\_check\_golze.py}):
\begin{center}\footnotesize
\begin{tabular}{p{7.6cm} rr}
\toprule
 & $V/g=0.707$ & $V/g=1.414$\\
\midrule
overlap$^2$ of $\Phi_c$ with the selected eigenstate of $H^c_{\rm imp}$ & $1.000000$ & $1.000000$\\
index of that eigenstate (0 = lowest; $H^c_{\rm imp}$ has $36$ states) & $3$ & $5$\\
$\braket{\Phi_c|H^c_{\rm imp}|\Phi_c}$ & $-22.04$ & $-22.04$\\
lowest eigenvalue of $H^c_{\rm imp}$ & $-24.32$ & $-30.19$\\
index of $\Phi_o$, $\Phi_v$ in their $H_{\rm imp}$ ($4$ states each) & $0$, $0$ & $0$, $0$\\
bath levels of $\lambda^c_c$ (eigenvalues of the $3\times3$ matrix) & $-2.14$, $-1.73$, $+1.73$ & $-5.07$, $-3.00$, $+3.00$\\
eigenvalues of ${}^c\Delta$ & $6.0\times10^{-6}$, $0.9979$, $1.0000$ & $2.5\times10^{-5}$, $0.9978$, $1.0000$\\
\bottomrule
\end{tabular}
\end{center}
That is: $\Phi_c$ is an exact eigenstate of $H^c_{\rm imp}$, as Eq.~\eqref{eq:eig} requires; it is not the lowest one (the fourth and the sixth of the $36$ eigenstates, with the
lowest eigenvalue $2.3$ and $8.1$ below it), whereas $\Phi_o$ and $\Phi_v$ are ground states. The dimensions are those of the half-filled sector $N_\uparrow=N_\downarrow=(1+B)/2$ of the fragment (physical orbital plus $B$ bath orbitals per spin):
for the core, $4$ orbitals per spin with $2$ electrons of each spin, $\binom42^2=36$; for $o$ and $v$, $2$ orbitals per spin with $1$ electron of each spin, $\binom21^2=4$.

\paragraph{How the size of $H^I_{\rm imp}$ is obtained.} The embedding problem of fragment $I$ has the physical orbital $c$ and $B_I$ bath orbitals $d_a$ per spin
(the bath orbitals include the ghosts: $B_c=3$ is one bath level plus two ghosts), i.e.\ $M=B_I+1$ orbitals per spin and $2M$ spin orbitals. The full Fock space has dimension $2^{2M}$.
The Hamiltonian $H^I_{\rm imp}$ conserves $N_\uparrow$ and $N_\downarrow$, and the code (\texttt{Fragment} in \texttt{gga\_cederbaum.py}) projects it on the sector with
$N_\uparrow=N_\downarrow=M/2$ (half filling, which requires $B_I$ odd). Spin up and spin down are independent in the counting, so
\begin{equation}
\dim H^I_{\rm imp}=\binom{M}{M/2}^2,\qquad M=B_I+1 .
\end{equation}
The operators are built in the full Fock space with sparse Jordan--Wigner matrices and then restricted to this sector (stored dense), and the dimension printed above is the shape of the projected
\texttt{Hloc}. The values are
\begin{center}\small
\begin{tabular}{c c c c}
\toprule
$B_I$ & $M$ & full Fock space $2^{2M}$ & sector $N_\uparrow=N_\downarrow=M/2$\\
\midrule
1 & 2 & $16$ & $4$\\
3 & 4 & $256$ & $36$\\
5 & 6 & $4096$ & $400$\\
\bottomrule
\end{tabular}
\end{center}
so $B_c=3$ and $B_c=5$ give $36$ and $400$ states for the core, and $4$ for each of $o$ and $v$ ($B_o=B_v=1$). The number of states grows much faster than the number of ghosts, but
the solution only needs the one eigenstate $\Phi_I$ that satisfies the constraint, not all of them. The index of the excited state is not a property of the model alone: it
depends on the multipliers recovered at the point, and these contain one direction that the stationarity conditions leave undetermined (Sec.~\emph{Absence of satellites} of \texttt{cederbaum.tex}).

%======================================================================
\section{Why it can be excited}

\begin{enumerate}
\item \textbf{Stationarity is not minimality.} Stationarity of $\mathcal L$ in $\Phi$ gives Eq.~\eqref{eq:eig} for \emph{some} eigenstate. It is the
\emph{constraint}, not the energy of $H_{\rm imp}$, that selects which one.
\item \textbf{The constraint fixes the bath occupations.} ${}^c\Delta=\mathbb I-\Delta^{bb}$ must equal $\rho^{cc}$, the local density matrix of the
quasiparticle state. With one bath electron per spin, this fixes \emph{which} bath orbital is filled: the one that corresponds to the empty
quasiparticle orbital. The multipliers $\lambda_I$ are then adjusted so that this state is stationary.
\item \textbf{That need not be the lowest state at the recovered multipliers.} At these multipliers the lowest state of $H^c_{\rm imp}$ would put the
bath electron in the lowest bath level ($-2.14$ and $-5.07$ at the two couplings) (and, with it, change the core occupation), which contradicts the constraint satisfied by $\Phi_c$. The state that
satisfies the constraint is a different, higher eigenstate.
\item \textbf{Non-convexity.} The constraint is quadratic in $\Phi$ and the energy functional is nonlinear in the quantities extracted from $\Phi$
(through $R_I=S^{-1}D$ and $\bar\Theta$), so the problem is not convex and a duality gap is possible: the minimum of $\mathcal L$ over $\Phi$ at the
optimal multipliers is not the constrained minimum. This is why the saddle-point assumption of the loop can fail.
\end{enumerate}
Points 1 and 2 follow directly from the equations; points 3 and 4 are my reading of the numerical result, not a proof.

%======================================================================
\section{Consequences for the impurity loop}

\paragraph{Lowest-state selection.} Started from the full direct-minimization state (\texttt{state0}), one loop iteration already moves away: the residual
$\|\rho^{II}-(1-\Delta^{bb}_I)\|$ is $1.41$ at both couplings (shift $\varepsilon=10^{-9}$), and after 300 iterations no mixing ($1$, $0.5$, $0.2$) has converged. The loop ends either in
the $B_c=1$-like state ($\nu_c=1.9959$, single pole at $-10.02$, $E$ higher than the direct minimum by $1.6\times10^{-4}$ and $6.4\times10^{-4}$) or in a state with $R_c=0$.

\paragraph{Maximum-overlap (MOM-like) selection.} In each embedding solve, the eigenstate with the largest overlap with the previous
$\Phi_I$ is taken instead of the lowest one (\texttt{follow=True}; degenerate subspaces are projected). One step from the direct-minimization
state then gives a residual of $6.1\times10^{-2}$ ($V/g=0.707$) and $2.2\times10^{-3}$ ($V/g=1.414$), instead of $1.41$, but it is still not a fixed point at $\varepsilon=10^{-9}$.
A filled ghost orbital has ${}^c\Delta$ eigenvalue $1.0000$, so $S_c\sim\sqrt\varepsilon$ and $\lambda^c$, obtained from $S_c^{-1}$, is ill conditioned
for small $\varepsilon$ (this is my reading of the mechanism). The followed state is not in a dense part of the spectrum here: its nearest neighbours are $1.1$ and $0.9$
($V/g=0.707$), $2.0$ and $0.14$ ($V/g=1.414$) away, apart from degenerate pairs at lower energy, which are projected.
After that the loop drifts to the $B_c=1$-like state (at $V/g=1.414$, mixing $0.2$ stays there, $E$ higher by $6.4\times10^{-4}$, pole $-10.01$) or to the trivial product state ($E=-22$, $R_c=0$).

\paragraph{Dependence on the shift} (300 iterations, mixings $1.0$, $0.5$, $0.2$, started from the direct-minimization state; ``none'' means no convergence to $10^{-10}$ in the residual):
\begin{center}\footnotesize
\begin{tabular}{rp{2.6cm}p{4.6cm}p{4.6cm}}
\toprule
shift $\varepsilon$ & lowest-state selection & max-overlap, $V/g=0.707$ & max-overlap, $V/g=1.414$\\
\midrule
$10^{-9}$ & none & none (ends in the $B_c=1$-like or the product state) & none\\
$10^{-8}$ & none & mixings $0.5$, $0.2$ converge to the product state ($R_c=0$) & same\\
$10^{-7}$ & none & none & mixings $0.5$, $0.2$ converge to the product state\\
$10^{-6}$ & none & \textbf{mixings $1.0$, $0.2$ converge} (residual $\le3\times10^{-12}$) to a nearby minimum, $E$ higher by $1.2\times10^{-4}$, pole $-10.017$ & mixings $0.5$, $0.2$ converge to the product state\\
$10^{-5}$ & none & none (stays at the $B_c=1$-like state, residual $0.7$--$1.4$) & \textbf{mixing $0.2$ converges} (residual $6\times10^{-13}$) to a nearby minimum, $E$ higher by $3.5\times10^{-4}$, pole $-9.970$\\
$10^{-4}$ & none & none (residual $4\times10^{-7}$--$2\times10^{-6}$ after 300 iterations, $E$ higher by $0.9$--$4\times10^{-6}$) & none\\
$10^{-3}$ & none & mixing $0.2$ converges to the product state & none (product state)\\
\bottomrule
\end{tabular}
\end{center}
The shift also changes the problem, so the direct-minimization energy used as reference changes with $\varepsilon$ (at $10^{-4}$ and $10^{-3}$ it is $-22.0412$ and $-22.0195$ at $V/g=0.707$).
At the converged nearby minima ($\varepsilon=10^{-6}$, $V/g=0.707$, mixing $0.2$, and $\varepsilon=10^{-5}$, $V/g=1.414$, mixing $0.2$) restarting the direct minimizer from the loop point lowers $E$ by
$4\times10^{-11}$ and $4\times10^{-11}$ (by $\le10^{-14}$ for mixing $1.0$ at $V/g=0.707$): the loop point is itself a constrained local minimum, with the same single pole in $G^h_{cc}$.

\paragraph{Conclusion.} The impurity loop reproduces a solution of this kind only with maximum-overlap state selection and a large regularization shift ($10^{-6}$ at $\rho/\Omega=0.5$, $10^{-5}$ at $\rho/\Omega=1.0$; the other shifts give no
convergence or the trivial product state); and then it finds a different, higher local minimum. The direct minimization needs no choice of eigenstate and is the more reliable solver.

\paragraph{Scope.} The numbers are for $V/g=0.707$ and $1.414$, $t_c=0.5$, $B_c=3$ (\texttt{excited\_check\_golze.py} and \texttt{shift\_sweep.py}; logs in \texttt{cederbaum\_data/logs}). Which mixing converges depends on the starting point and on the one direct-minimization solution found
(best of $5$ random starts, $30$ at $\varepsilon=10^{-3}$, $V/g=0.707$), so the table is a sample, not a map of all behaviour. The explanation in Sec.~4, points 3--4, is an interpretation of
those numbers.

\end{document}
